\documentclass{article}
\usepackage[margin=1in]{geometry}
\usepackage{mathtools, amsfonts, amsthm, xcolor, listings}

\newcommand{\R}{\mathbb{R}}
\newcommand{\N}{\mathbb{N}}
\newcommand{\I}{\mathbb{I}}
\newcommand{\Q}{\mathbb{Q}}
\newcommand{\Z}{\mathbb{Z}}

\newcommand{\set}[1]{\left\{#1\right\}}
\newcommand{\ang}[1]{\left\langle #1 \right\rangle}
\newcommand{\inv}{^{-1}}

\newcommand{\normsub}{\lhd}

\newcommand{\reflection}[5]{
    \begin{itemize}
    \item Points attempted: #1
    \item Self-assessed score: #2/#1.
    \item Time spent: #3 minutes.
    \item Difficulty: #4/10.
    \item Questions/Feedback: #5.
    \end{itemize}
}


\title{HW09}
\author{Sam Ly}

\begin{document}
\maketitle

\section*{Chapter 18}
https://github.com/samlyme/Assignments/tree/cs4080/ch18/main

\begin{enumerate}
    \item {
        We could reduce our binary operators even further than we did here. Which other instructions can you eliminate, and how would the compiler cope with their absence?

        The two operations that could potentially be removed are the OP NIL, TRUE, and FALSE,
        and the either OP GREATER or OP LESS.

        We can remove all the data creation ops simply because we already have the 
        OP CONSTANT. Then, for OP GREATOR, the compiler can just reorder the elements 
        and use an OP LESS and vice versa.
    }
    \item {
        Conversely, we can improve the speed of our bytecode VM by adding more specific instructions that correspond to higher-level operations. What instructions would you define to speed up the kind of user code we added support for in this chapter?

        We can add more boolean operations like NOT EQUAL and GREATER THAN OR EQUAL.

    }
\end{enumerate}

\end{document}
